\section{OCUDU mit PDSCH-Traffic starten und die Generierung beobachten}

Es gibt zwei Wege: die komplette gNB im Testmodus mit Dummy-RU (kein Core, kein Funk, aber echter PDSCH-Traffic
für ein Phantom-UE), oder die Unit-Tests der Codebook-Generierung im Debugger.

\subsection{Bauen}

\begin{lstlisting}[language=bash]
cmake --build cmake-build-debug --target gnb precoding_codebook_generator_test precoding_codebooks_test -j4
\end{lstlisting}

\subsection{Konfiguration: Testmodus mit 4 Ports}

Vorlage ist \file{configs/testmode.yml}. Für Type-I-Codebooks sind mindestens 4 DL-Antennen nötig, sonst wird
nur das 2-Port- oder 1-Port-Codebook erzeugt.

\begin{lstlisting}[language=yaml,title={configs/testmode\_4x4.yml (Beispiel)}]
cu_cp:
  amf:
    no_core: true

ru_dummy:
  dl_processing_delay: 1

cell_cfg:
  dl_arfcn: 632628
  band: 78
  channel_bandwidth_MHz: 20
  common_scs: 30
  plmn: "00101"
  tac: 7
  pci: 1
  nof_antennas_dl: 4
  nof_antennas_ul: 4
  # Optional: Topologie mit Beam-Raster. Ohne Angabe wird four_ports abgeleitet.
  # tx_ant_topology: single_panel_two_one

test_mode:
  test_ue:
    rnti: 0x44
    nof_ues: 1
    pdsch_active: true
    pusch_active: false
    cqi: 15
    ri: 2          # 1..4, bestimmt die Layer-Anzahl
    i_1_1: 1       # PMI-Felder, die das Phantom-UE "meldet"
    i_1_3: 0
    i_2: 1

log:
  filename: /tmp/gnb.log
  all_level: warning
  mac_level: debug
  phy_level: debug
\end{lstlisting}

Das Phantom-UE meldet seine CSI aus genau diesen Werten. Das Packen passiert in
\file{lib/du/du\_high/test\_mode/mac\_test\_mode\_helpers.cpp}:

\begin{lstlisting}[language=C++,title={lib/du/du\_high/test\_mode/mac\_test\_mode\_helpers.cpp, fill\_pmi\_bits()}]
const auto* single_panel = std::get_if<pmi_codebook_typeI_single_panel>(&pmi_codebook);
if (single_panel == nullptr) {
  return;
}

// The fields that the panel topology and the rank do not report have a zero bit-width.
const pmi_typeI_single_panel_param_sizes sizes =
    get_pmi_sizes_typeI_single_panel(get_single_panel_info(single_panel->n1_n2), rank);
packed.push_back(test_ue_cfg.i_1_1, sizes.i_1_1);
packed.push_back(test_ue_cfg.i_1_2.value_or(0), sizes.i_1_2);
packed.push_back(test_ue_cfg.i_1_3.value_or(0), sizes.i_1_3);
packed.push_back(test_ue_cfg.i_2, sizes.i_2);
\end{lstlisting}

Starten:
\begin{lstlisting}[language=bash]
./cmake-build-debug/apps/gnb/gnb -c configs/testmode_4x4.yml
\end{lstlisting}

Die Topologie-Namen liefert \code{./cmake-build-debug/apps/gnb/gnb cell\_cfg -h}. Namen mit Präfix
\code{single\_panel\_} haben ein Beam-Raster.

\subsection{Breakpoints für die Generierung}

Die Codebook-Erzeugung loggt nichts und läuft genau einmal beim Start, bevor Traffic fließt.
\begin{itemize}
  \item \file{lib/fapi\_adaptor/precoding\_codebook\_generator.cpp:184}, die Zeile
    \code{repo\_builder.add(pm\_index, make\_composite(...))} in der vierfachen Schleife.
  \item \file{lib/ran/precoding/type1\_sp/precoding\_codebooks.cpp}, Einstieg von \code{make\_type1\_sp\_mode1()}
    und der \code{switch (nof\_layers)}.
  \item \file{lib/du/du\_low/o\_du\_low\_factory.cpp:64} und \file{lib/du/du\_high/o\_du\_high\_factory.cpp:105},
    die beiden Aufrufer.
\end{itemize}

Isolierter geht es über die Unit-Tests:
\begin{lstlisting}[language=bash]
./cmake-build-debug/tests/unittests/fapi_adaptor/precoding_codebook_generator_test --gtest_filter='*TypeI*'
./cmake-build-debug/tests/unittests/ran/precoding_codebooks_test
\end{lstlisting}

\subsection{Die Nutzung im laufenden Traffic}

Mit \code{mac\_level: debug} schreibt der Scheduler pro PDSCH-Grant die gewählte PMI ins Log
(\file{lib/scheduler/logging/scheduler\_result\_logger.cpp:627}, Format \code{ri=<layers>} gefolgt von den
PMI-Feldern). Mit \code{phy\_level: debug} protokolliert der PDSCH-Prozessor die angewendete
Precoding-Konfiguration.

\begin{lstlisting}[language=bash]
grep -E "PDSCH.*ri=" /tmp/gnb.log | head
\end{lstlisting}

Den Weg von der PMI im Scheduler zum Index im Repository zeigt \file{lib/fapi\_adaptor/precoding\_codebook\_mapper.cpp}.
Zur Laufzeit wird nichts neu gerechnet: Der PHY-Adaptor holt in
\file{lib/fapi\_adaptor/phy/p7/pdu\_translators/precoding.h:19} mit \code{get\_precoding\_config()} das fertige
Composite aus dem Repository.
